\rok{Јануар 2026.}{B}

Први поправни тест из Алгоритама и структура података

Први практични тест одржан 29.01.2026.

\zadatak{1}{Insertion Sort}
Написати методу \texttt{sort} Insertion Sort алгоритма за сортирање целобројног низа дужине $n$.

\textbf{Додатне напомене за решавање задатка:}
\begin{itemize}
	\item Улаз је несортиран низ целих бројева.
	\item Излаз је сортирани низ целих бројева.
	\item У template-у се налази класа \texttt{RandomArrayGenerator} коју треба искористити за генерисање низа случајних бројева.
	\item У оквиру класе \texttt{InsertionSort} написати имплементацију за методу:
	\begin{Verbatim}[fontsize=\small]
private static void sort(long[] arr)
	\end{Verbatim}
	\item Обезбедити код у Main методи за тестирање алгоритма.
\end{itemize}

\zadatak{2}{Проналажење "Балансне Тачке" листе}
Потребно је имплементирати методу која проналази вредност првог чвора у листи који представља "балансну тачку". Чвор је балансна тачка ако је сума свих елемената са његове леве стране једнака суми свих елемената са његове десне стране.

\textbf{Додатни захтеви за решавање задатка:}
\begin{itemize}
	\item Ако такав чвор не постоји, метода треба да врати \texttt{-1}.
	\item Задатак се мора решити у линеарном времену $O(n)$.
	\item Сматрати да је сума "празне стране" (код првог или последњег чвора) једнака 0.
	\item Захтеван потпис методе:
	\begin{Verbatim}[fontsize=\small]
public int findBalancePoint()
	\end{Verbatim}
\end{itemize}

\textbf{Пример:}
\begin{itemize}
	\item \textbf{Улаз 1:} \texttt{4 -> 3 -> 6 -> 7}
	\item \textbf{Излаз 1:} \texttt{6}
	\item \textbf{Улаз 2:} \texttt{2 -> 3 -> 4 -> 4}
	\item \textbf{Излаз 2:} \texttt{-1}
\end{itemize}

\zadatak{3}{Чишћење мрежног бафера (Network Buffer Scrubber)}
Замислите мрежни бафер (Queue) који прима пакете података (позитивни бројеви) који чекају на обраду. Повремено у систем улази "Antisignal" или "шум" (негативан број).

\textbf{Правила интеракције:}
\begin{enumerate}
	\item Када стигне Antisignal (негативан број), он напада пакет који је тренутно први у реду (Front).
	\item Судар:
	\begin{enumerate}[a.]
		\item Ако је пакет на почетку реда мањи од апсолутне вредности antisignala, пакет је уништен, а antisignal наставља да уништава следећи пакет у реду са преосталом јачином.
		\item Ако је пакет исти као апсолутна вредност antisignala, оба нестају.
		\item Ако је пакет већи од апсолутне вредности antisignala, пакет преживљава, али његова вредност се смањује за јачину antisignala. Квака: пошто је пакет "обрађен" овим сударом, он се пребацује на крај реда (Back) како би уступио место другим пакетима.
	\end{enumerate}
	\item Ако antisignal уништи све пакете у реду и још увек има преосталу јачину, он сам постаје пакет (негативне вредности) и улази у ред на крај.
\end{enumerate}

\textbf{Додатни захтеви за решавање задатка:}
\begin{itemize}
	\item Временска сложеност: $O(n)$.
	\item Захтевани потпис методе:
	\begin{Verbatim}[fontsize=\small]
public int[] ScrutinizeBuffer(int[] data)
	\end{Verbatim}
\end{itemize}

\textbf{Пример:}
\begin{itemize}
	\item \textbf{Улаз 1} (несортирани низ целих бројева): \texttt{[10, 5, 20, -25]}
	\item \textbf{Кораци:}
	\begin{itemize}
		\item \texttt{10, 5, 20} улазе редом: Queue: \texttt{[10, 5, 20]}
		\item Стиже \texttt{-25} (noise = 25).
		\item Напада \texttt{10} (front): \texttt{10} је мањи, \texttt{10} нестаје. Noise остаје \texttt{15}.
		\item Напада \texttt{5} (front): \texttt{5} је мањи, \texttt{5} нестаје. Noise остаје \texttt{10}.
		\item Напада \texttt{20} (front): \texttt{20} је већи.
	\end{itemize}
	\item \textbf{Излаз 1:} \texttt{[10]}
\end{itemize}

\sablon{Шаблон првог поправног теста јануар 2026. група Б}{2025/Januar/GrupaB/AISP2025_PopravniTest1_Test1_GrupaB.linq}
